\chapter{Code}
\section{Autonomous Delivery Vehicle: Arduino Code}
The following code snippet is used to control the autonomous delivery vehicle using a NodeMCU board and Blynk. The code manages servo motors and DC motor control for directional movement based on joystick inputs from the Blynk app.

\begin{tcblisting}{listing only,
  colback=gray!10!white, 
  breakable, 
  boxsep=0pt,
  top=1mm,
  bottom=1mm,
  left=1mm,
  right=1mm,
  listing options={
    language=C++,
    basicstyle=\small\ttfamily, 
    keywordstyle=\color{blue}, 
    commentstyle=\color{gray},
    backgroundcolor=\color{gray!25},
    escapechar=!,
    morekeywords={BLYNK_WRITE, BlynkEdgent, analogWrite, digitalWrite, pinMode, attach}
  }
}
//
// NodeMCU Blynk-Controlled Autonomous Vehicle
//

#define BLYNK_TEMPLATE_ID "TMPL3UJinHw1e"
#define BLYNK_DEVICE_NAME "AutonomousDeliveryVehicle"

#define BLYNK_FIRMWARE_VERSION        "0.1.0"
#define BLYNK_PRINT Serial
#define USE_NODE_MCU_BOARD

#include "BlynkEdgent.h"
#include <Servo.h>

#define servo1 D2
#define servo2 D5
#define servo3 D1
#define servo4 D5

#define ENA D6
#define IN1 D7
#define IN2 D8
#define IN3 D3
#define IN4 D0
#define ENB D4

Servo mservo1, mservo2, mservo3, mservo4;
int x = 50;
int y = 50;
int Speed = 255;

BLYNK_WRITE(V0) {
  int s0 = param.asInt(); 
  mservo1.write(s0);
}
BLYNK_WRITE(V1) {
  int s0 = param.asInt(); 
  mservo2.write(s0);
}
BLYNK_WRITE(V2) {
  int s0 = param.asInt(); 
  mservo3.write(s0);
}
BLYNK_WRITE(V3) {
  int s0 = param.asInt(); 
  mservo4.write(s0);
}
BLYNK_WRITE(V4) {
  x = param[0].asInt();
}
BLYNK_WRITE(V5) {
  y = param[0].asInt();
}

void smartcar() {
  if (y > 70) {
    carForward();
    Serial.println("carForward");
  } else if (y < 30) {
    carBackward();
    Serial.println("carBackward");
  } else if (x < 30) {
    carLeft();
    Serial.println("carLeft");
  } else if (x > 70) {
    carRight();
    Serial.println("carRight");
  } else {
    carStop();
    Serial.println("carstop");
  }
}

void carForward() {
  analogWrite(ENA, Speed);
  analogWrite(ENB, Speed);
  digitalWrite(IN1, HIGH);
  digitalWrite(IN2, LOW);
  digitalWrite(IN3, LOW);
  digitalWrite(IN4, HIGH);
}

void carBackward() {
  analogWrite(ENA, Speed);
  analogWrite(ENB, Speed);
  digitalWrite(IN1, LOW);
  digitalWrite(IN2, HIGH);
  digitalWrite(IN3, HIGH);
  digitalWrite(IN4, LOW);
}

void carLeft() {
  analogWrite(ENA, 0);
  analogWrite(ENB, Speed);
  digitalWrite(IN1, LOW);
  digitalWrite(IN2, HIGH);
  digitalWrite(IN3, LOW);
  digitalWrite(IN4, HIGH);
}

void carRight() {
  analogWrite(ENA, Speed);
  analogWrite(ENB, 0);
  digitalWrite(IN1, HIGH);
  digitalWrite(IN2, LOW);
  digitalWrite(IN3, HIGH);
  digitalWrite(IN4, LOW);
}

void carStop() {
  digitalWrite(IN1, LOW);
  digitalWrite(IN2, LOW);
  digitalWrite(IN3, LOW);
  digitalWrite(IN4, LOW);
}

void setup() {
  Serial.begin(9600);
  mservo4.attach(servo1);
  mservo2.attach(servo2); 
  mservo3.attach(servo3); 
  mservo4.attach(servo4);

  pinMode(ENA, OUTPUT);
  pinMode(IN1, OUTPUT);
  pinMode(IN2, OUTPUT);
  pinMode(IN3, OUTPUT);
  pinMode(IN4, OUTPUT);
  pinMode(ENB, OUTPUT);
  BlynkEdgent.begin();
  delay(1000); 
}

void loop() {
  BlynkEdgent.run();
  smartcar();
}
\end{tcblisting}
